% Part: computability
% Chapter: recursive-functions
% Section: normal-form

\documentclass[../../../include/open-logic-section]{subfiles}

\begin{document}

\olfileid[hi]{cmp}{rec}{nft}
\olsection{प्रसामान्य-रूप प्रमेय}

\begin{thm}[क्लिनी (Kleene) का प्रसामान्य-रूप प्रमेय]
\ollabel{thm:kleene-nf}
एक आदिम पुनरावर्ती संबंध $T(e, x, s)$ और एक आदिम पुनरावर्ती फलन $U(s)$
हैं जिनका निम्न गुण है: यदि $f$ कोई आंशिक पुनरावर्ती फलन है, तो किसी~$e$
के लिए
\[
f(x) \simeq U(\umin{s}{T(e, x, s)})
\]
प्रत्येक $x$ के लिए मान्य है।
\end{thm}

\begin{explain}
प्रसामान्य-रूप प्रमेय का प्रमाण विस्तृत है, किंतु मूल विचार सरल है। हर
आंशिक पुनरावर्ती फलन का कोई \emph{अनुक्रमांक}~$e$ होता है, जिसे सहज रूप
से उसके प्रोग्राम अथवा परिभाषा को कूटबद्ध करने वाली संख्या समझें। यदि
$f(x)
\fdefined$, तो संगणन को व्यवस्थित रूप से दर्ज करके किसी संख्या~$s$
से कूटबद्ध किया जा सकता है; और इस तथ्य को कि $s$,~$f$ के निवेश~$x$ वाले
संगणन को कूटबद्ध करता है, केवल $x$ और परिभाषा~$e$ का उपयोग करके आदिम
पुनरावर्ती ढंग से जाँचा जा सकता है। फलतः संबंध~$T$---``अनुक्रमांक~$e$
वाले फलन का निवेश~$x$ के लिए कोई संगणन है और $s$ उस संगणन को कूटबद्ध
करता है''---आदिम पुनरावर्ती है। संगणन का पूरा अभिलेख~$s$ दिए होने पर
~$s$ का ``परिणाम''~$f(x)$ का मान है और उसे भी~$s$ से आदिम पुनरावर्ती
ढंग से पाया जा सकता है।

प्रसामान्य-रूप प्रमेय दिखाता है कि किसी भी आंशिक पुनरावर्ती फलन की
परिभाषा के लिए केवल एक अपरिबद्ध खोज चाहिए। मूलतः हम सभी संख्याओं में तब
तक खोज सकते हैं जब तक ऐसी संख्या न मिले जो अनुक्रमांक~$e$ वाले फलन के
निवेश~$x$ के किसी संगणन को कूटबद्ध करती हो। हम संख्याओं~$e$ को आंशिक
पुनरावर्ती फलनों के ``नाम'' के रूप में उपयोग कर सकते हैं और प्रमेय के
समीकरण से परिभाषित फलन के लिए $\cfind{e}$ लिख सकते हैं; वह फलन~$f$ है। ध्यान दें कि
किसी आंशिक पुनरावर्ती फलन के एक से अधिक अनुक्रमांक हो सकते हैं---वास्तव
में प्रत्येक आंशिक पुनरावर्ती फलन के अनंततः अनेक अनुक्रमांक होते हैं।
\end{explain}
\end{document}
